\documentclass{article}
\usepackage{pathfinder-readable}

\title{When Peer Agreement Measures a Shared Error}
\pathfinderpapers{2609.01595}{Mechanism Design for Alignment and Control}%
  {2607.27967}{MARS-RA: Rank Aggregation for Credit Assignment via Multimodal
  Comparisons in Embodied Multi-Agent Cooperation}

\begin{document}
\maketitle

\begin{abstract}
This account connects a mechanism for eliciting AI agents' private types in Q with a method for
assigning team credit from image comparisons in P. The thread asked whether truthful peer reports
and repeated comparisons can certify the contribution the designer cares about. The peers found a
shared-error example in which reports agree and comparisons concentrate, yet neither identifies the
target. Independent observations or audits give limited ways to measure the error, subject to
separate incentives for truthful reporting. The verifier paused the thread because this known
identification limit does not yet yield a new result for P's setting or evidence on its benchmark.
\end{abstract}

\section{Introduction}

Bergemann, Koh and Morris study how a designer can induce AI agents to reveal what they know and do
what they are asked to do \cite{Q}. An agent's preferences, abilities and information may be
private. In their multi-agent example, a reward scores each agent's reported type against the types
reported by peers. Under a condition that different types predict different peer types, this helps
make truthful reporting an equilibrium. Further parts of the reward induce the desired action.

Wang and co-authors study how to assign credit to agents cooperating in embodied tasks \cite{P}.
Their MARS-RA method asks a large multimodal model to compare agents' current images, fits scores to
those pairwise choices using a Bradley--Terry model, and uses the scores to shape rewards during
training. Its convergence claim concerns the model's underlying preference. Its Shapley-value
interpretation assumes that this preference already agrees with Shapley values, a measure of
contribution across possible teams.

The connection pursued was whether Q's peer scoring could support P's claim that aggregated
comparisons measure useful contribution. Three questions emerged: whether an agent has reason to
report what it sees, whether comparisons recover the comparator's preference, and whether that
preference tracks the intended contribution. These are different claims.

\section{What was done}

The peers first read Q's peer-discipline result beside P's rank-aggregation result. Ada checked Q's
separation condition and the reward calculation; Emmy checked what P's fitted score is meant to
recover. They found that Q assumes a specified map from each private type to a distribution of peer
types. P assumes a stable latent preference for its comparison model, then separately assumes
equality with Shapley values for that interpretation. Neither paper derives target accuracy from
peer agreement alone.

They next tested a shared observation error. They treated an observed bit as Q's private type and
checked the score against a truthful peer. They also considered repeated comparisons of the same
visual scene in P. Ada checked an independent audit and the existence of P's score estimate; Emmy
checked independent observation groups and the timing of the two mechanisms. Each peer cross-checked
the other's calculations. The worked examples and qualifications are in
\url{../ada/peer_scoring_note.md} and \url{../emmy/shared_noise.md}.

\section{Results}

Q's condition separates types by their predictions about peers. Let \(\beta_i[t]\) be the
distribution of other agents' types predicted by agent \(i\) when its type is \(t\). Against
truthful peers, Q's quadratic score gives type \(t\) an expected advantage over a false report \(u\)
of
\[
\Lambda\lVert\beta_i[t]-\beta_i[u]\rVert_2^2,
\]
where \(\Lambda>0\) scales the reward. When the two predicted distributions differ, this advantage
is positive. Q also cancels the agent's utility on the intended action support and penalises actions
outside it. Thus its result supports a truthful, obedient equilibrium under its stated conditions;
the peer score alone is not the obedience argument.

The shared-error example shows the limit of that result as a measurement claim. Let \(T\) be a fair
binary target, \(N\) an independent binary error with \(\Pr(N=1)=p\), and let every peer observe
\(X_i=T\oplus N\), where \(\oplus\) means addition modulo two. Each peer's observation predicts
every other peer's observation exactly, so Q's separation condition holds for the two possible
observed types. Yet the full collection of peer observations is just one fair bit for every value of
\(p\). When \(p=0\), it gives the target; when \(p=1/2\), it gives no information; when \(p=1\), it
gives the reverse. Even with a known error rate below one half, adding peers who share this error
does not reduce the error on a task. Ada and Emmy checked the conditional distributions and the
resulting score gap in their notes.

The same distinction appears in P's comparisons. Let \(d\) be the intended difference in
contribution between two agents, and \(B\) a shift shared by all comparisons of one scene. If
repeated comparison outcomes are independent conditional on \(B\), with probability \(\sigma(d+B)\)
of preferring the first agent, where \(\sigma(z)=1/(1+\exp(-z))\), then a two-agent Bradley--Terry
fit converges to \(d+B\). More calls refine the scene-conditioned preference. They do not separate
\(d\) from \(B\): raising one and lowering the other by the same amount leaves the comparison law
unchanged. This is a model-based example, not an empirical finding about P's benchmark.

The peers derived two restricted ways to check target accuracy. If \(G\) observation groups have
independent binary errors with the same probability \(q<1/2\), and each group supplies one truthful
report, a majority vote has error at most
\[
\exp\{-2G(1/2-q)^2\}.
\]
Here \(G\) counts independent groups, not repeated calls within one group. Alternatively, let an
independent audit label be \(A=T\oplus E\), where \(E\) has known error probability \(q<1/2\). Under
truthful reports, the report error rate \(p\) follows from
\[
\Pr(A=X)=1-q-p(1-2q),\qquad
p=\frac{1-q-\Pr(A=X)}{1-2q}.
\]
Ada derived the audit identity and Emmy checked it. The majority bound concerns a binary target;
neither calculation establishes cardinal Shapley values. An audit used only to assess reports does
not make truth-telling attractive.

\section{Objections and limits}

The first objection was about timing. Q's agents report types before later private signals arrive.
P's model compares images after observation. A direct transfer of Q's theorem would need a new game
in which agents can report those later observations and receive suitable payments. If P's agents
cannot change the images or the comparison calls, the immediate question is whether the comparator
is calibrated to contribution, not whether agents report truthfully.

The second objection concerned equilibrium selection. In the binary example with two agents, zero
utility and one action, a peer receives score \(1\) for a matching report and \(-1\) otherwise.
Truthful reporting is a strict equilibrium, but both agents flipping their observations is another
strict equilibrium. Ada checked this case. Her audit-payment condition selects truthful reporting
only in that small game; it does not cover Q's general action problem. Earlier work cited in the
ledger discusses relabelling equilibria and informed truthfulness \cite{relabel,informed}.

The third objection concerned P's stated finite-sample guarantee. P defines an unregularised
Bradley--Terry estimate but its proof sketch refers to regularised-estimator analysis. A connected
comparison graph does not by itself ensure a finite unregularised estimate: with two agents and
every comparison favouring the same one, the likelihood keeps increasing as their fitted score gap
grows. Ada's two-agent check establishes this exception to the printed assumptions. An existence
condition or specified regulariser would address that issue, without proving calibration to
contribution.

Finally, P's potential-based reward shaping has a different role from Q's direct peer payment. With
a fixed initial state and zero terminal potential, its discounted shaping rewards sum to a constant
set by the initial potential. They can change training feedback without changing the exact ranking
of complete trajectories under those conditions. They do not supply Q's reporting and obedience
incentives. The ledger also cites work on measurement integrity, conditional independence in peer
elicitation, and random effects in Bradley--Terry comparisons \cite{integrity,independence,random}.
The peers did not independently assess those sources, and did not claim the general shared-error
warning as new.

\section{Why the thread stopped}

The verifier accepted the shared-error counterexample but judged it a known identification limit. It
did not establish a new result connecting Q's pre-signal reporting mechanism to P's post-observation
comparisons. A stronger claim would need a declared contribution target and independent audits or
observation groups, backed by a new theorem or benchmark evidence. The peers had neither a general
certificate for P's changing teams nor a test of the proposed dependence effect in its benchmark.

The smallest restart step is a controlled test against one declared target in P's benchmark. At a
fixed comparison budget, compare repeated judgements of one scene with judgements from independently
generated observation groups, using held-out dense-reward labels and uncertainty measured by scene.
That would test the practical size of shared error for that target. It would leave any strategic
reporting mechanism for a separately specified game.

\bibliographystyle{plainurl}
\bibliography{references}
\end{document}
